\fichatitulo{27 · REVISÃO}{Preprint v1 · 2024}{Revisão abrangente de RAG}
{\small Gupta et al. \textit{A Comprehensive Survey of Retrieval-Augmented Generation (RAG): Evolution, Current Landscape and Future Directions}. Preprint v1 · 2024. \href{https://arxiv.org/pdf/2410.12837v1}{Fonte consultada}.\par}

\campo{Problema e objetivo}
Quais componentes e desafios estruturam as aplicações de RAG?

\campo{Principal contribuição · síntese interpretativa}
Reúne arquitetura, evolução, aplicações e desafios de sistemas que integram recuperação e geração.

\campo{Método / natureza da evidência}
Revisão narrativa; o recorte consultado concentra-se na organização e nas conclusões.

\campo{Resultado ou argumento principal}
Destaca problemas de escala, custos e qualidade dos dados, além de vieses e questões éticas.

\campo{Excertos-chave · seleção literal}
\begin{itemize}
\excerto{1}{scalability, bias, and ethical concerns}{Resumo.}
\excerto{2}{high computational overhead}{Conclusão.}
\end{itemize}

\campo{Limitações e análise crítica}
Análise da ficha: enumeração de aplicações e benefícios não demonstra eficácia em todos os domínios. Esta leitura focal não verifica individualmente todos os estudos citados.

\campo{Aproveitamento na dissertação · proposta}
Fundamentação: mapear dimensões do problema; sustentar comparações quantitativas com os experimentos originais.

\registro{Fonte consultada nesta preparação: Resumo e seção 7 (conclusão). Chave: \texttt{gupta2024comprehensiveRag}.}
